\chapter{胶子 PDF：两条主线的交汇之处}

胶子分布是领域困难的集中地：伴随表示算符、更大的统计噪声、单态道中
与夸克的电荷共轭混合——以及全球拟合中最大的理论不确定度。本章五篇
论文为库内的胶子 PDF 纲领提供了理论，并使梯度流与大动量两条主线
合流。

% ----------------------------------------------------------------------
\section{胶子准 PDF 的匹配}
\papercard{J.-H.\ Zhang, X.\ Ji, A.\ Sch\"afer, W.\ Wang, S.\ Zhao}
{在大动量有效理论中获取胶子部分子分布}
{Phys.\ Rev.\ Lett.\ \textbf{122}, 142001 (2019)}
{arXiv:1808.10824 | DOI: 10.1103/PhysRevLett.122.142001}
\whyselected{库内每个胶子准 PDF 计算的理论引擎（约 $\sim$9/99）。}
\physics{借助通过 Yang--Mills 拉氏量耦合的伴随辅助重夸克场作为规范
不变调节器，论文证明胶子准 PDF 算符的全部幂发散都定域于 Wilson 线
自能并可由质量重整化消除；随后构造了独立可乘重整化的关联算符——四个
非极化、三个螺旋度——计算其到光锥胶子 PDF 的单圈匹配，并给出胶子道的
LaMET 因子化公式。该分类还澄清了领头 twist 下哪些组合不含夸克单态
混合。}
\impact{库内 Fan 等人的胶子准 PDF 快报与 Lin 组的 LaMET 胶子系列实现的
正是这一框架；混合重整化论文把其减除改造到有限流算符上。}
\cardend

% ----------------------------------------------------------------------
\section{短距离下的胶子伪分布}
\papercard{I.\ Balitsky, W.\ Morris, A.\ Radyushkin}
{短距离处的胶子伪分布：前向情形}
{Phys.\ Lett.\ B \textbf{808}, 135621 (2020)}
{arXiv:1910.13963 | DOI: 10.1016/j.physletb.2020.135621}
\whyselected{HadStruc 整条胶子线的理论枢纽（约 $\sim$10/99）：库内每篇
胶子伪 ITD 提取都引用其匹配核。}
\physics{论文对双局域双胶子算符分类，分离出进入前向矩阵元的 twist-2
不变振幅；用背景场方法与热核技术在显式规范不变形式下计算胶子伪 ITD
的短距离展开，
\begin{equation}
  \mathcal{M}_g(\nu,z^2)\big|_{|z^2|\ \mathrm{small}}
  = C_g(\nu,z^2,\mu)\otimes g(x,\mu) + \mathcal{O}(z^2),
\end{equation}
把紫外对数与 DGLAP 演化分离，并给出约化 ITD 的匹配关系。姊妹篇
（PRD \textbf{105}, 014008 (2022), arXiv:2111.06797，库内亦广泛引用）
提供全部投影与夸克混合项的完整表达式。}
\impact{Khan 等、Egerer 等、Salas-Chavira 等、Fan--Lin 及物理点/连续
极限提取都经由这些核转换矩阵元——之于胶子伪 PDF，它们正如 Xiong 等
之于非单态夸克道。}
\cardend

% ----------------------------------------------------------------------
\section{涂抹准 PDF：流与部分子的相遇}
\papercard{C.\ Monahan, K.\ Orginos}{准部分子分布与梯度流}
{JHEP \textbf{03} (2017) 116}
{arXiv:1612.01584 | DOI: 10.1007/JHEP03(2017)116}
\whyselected{本库原文之一，也是第三章与第四章之间的显式桥梁（其小生境
内约 $\sim$5/99，在涂抹算符支线内普遍）：它把流的有限性定理引入部分子
物理。}
\physics{把式~(\ref{eq:quasipdf}) 中裸 Wilson 线算符替换为由流场
$B_\mu(t)$、$\chi(t)$ 在小流时间 $t$ 构造的版本：既然流乘积紫外有限
（式~\ref{eq:flow}），矩阵元具有光滑连续极限而无需手工减除线性发散。
配合 ringed fermion，波函数重整化保持可乘；小 $t$ 展开把涂抹准 PDF 的矩
与光锥矩以可计算系数相联系，系数服从类 DGLAP 方程组。Monahan 的后续
单圈分析（PRD \textbf{97}, 054507 (2018)，库内）提供匹配核并确认涂抹
不改变红外结构。}
\impact{这是库内《PDF 的梯度流：首次应用于介子》与 Shindler 流化
twist-2 矩提案的概念源头：一旦算符被流涂抹，第五章的重整化战役就简化
为普通匹配——代价是新增一个标度 $t$，其系统处理由 NNLO 双线性论文
提供。}
\cardend

% ----------------------------------------------------------------------
\section{自重整化（self-renormalization）}
\papercard{Y.-K.\ Huo et al.\ (Lattice Parton Collaboration)}
{准光锥关联函数在格点上的自重整化}
{Nucl.\ Phys.\ B \textbf{969}, 115443 (2021)}
{arXiv:2103.02965 | DOI: 10.1016/j.nuclphysb.2021.115443}
\whyselected{本库原文之一，LP3 纲领的主力方案（约 $\sim$6/99）：被
横率、分布振幅与 TMD 论文使用。}
\physics{利用发散自身的结构：把重整化矩阵元写成
\begin{equation}
  Q_R(z,P^z;\mu)=Z^{-1}(z,a,m_0)\;Q_{\rm lat}(z,P^z,a),
  \qquad Z(z,a)\propto e^{\,-c\,z/a}\times(\text{对数}),
\end{equation}
并通过跨格距拟合 $Q_{\rm lat}$ 本身同时确定 $c$ 与残余对数因子——无需
外部顶点计算、不依赖 Landau 规范、不需要辅助场。论文验证了所提取 $Z$
跨强子动量与道的普适性，并记录了 RI/MOM 风格顶点在长距离处可观的非
微扰污染，从而支持"只在短距离使用"的原则。}
\impact{在本库内，这是若干旗舰结果背后的方案（连续极限下的核子横率；
π/K 分布振幅；TMD 波函数）；其诊断直接催生了下文的混合方案设计。}
\cardend

% ----------------------------------------------------------------------
\section{混合重整化}
\papercard{X.\ Ji, Y.-S.\ Liu, A.\ Sch\"afer, W.\ Wang, Y.-B.\ Yang,
J.-H.\ Zhang, Y.\ Zhao}
{大动量有效理论中准光锥关联函数的混合重整化方案}
{Nucl.\ Phys.\ B \textbf{964}, 115311 (2021)}
{arXiv:2008.03886 | DOI: 10.1016/j.nuclphysb.2021.115311}
\whyselected{本库原文之一；LP3 纲领的现行标准方案（约 $\sim$10/99）：
十年鏖战后的务实综合。}
\physics{沿分离距离轴分区处理。对 $z<z_c$ 采用短距离比值方案（或
RI/MOM），此处微扰可信、无非微扰污染；对 $z>z_c$ 通过辅助场质量项
减除线性发散的 Wilson 线自能并与短距区连续衔接，再转换到
$\overline{\rm MS}$。方案携带一个参数 $m_0^2$（等价地 $z_c$），具有
内禀 $\mathcal{O}(a)$ 歧义，文章给出了确定程序——跨格距拟合，或经
cusp 反常量纲固定。配合动量空间 $\ln(p_z/\mu)$ 对数的重求和
（Su 等，NPB \textbf{991}, 116201 (2023), arXiv:2209.01236，库内）
以及 renormalon 自洽的功率修正（Zhang 等，PLB \textbf{844}, 138081
(2023)），它定义了现代流水线。}
\impact{本库后期每一项 LP3 成果——胶子 PDF、Boer--Mulders、重介子
LCDA、TMD 软函数——都运行在此方案或其自重整化变体上；这两篇论文合计
是语料库后半段被引最多的方法论搭档。}
\cardend
